%%=============================================================================
%% Hoofdstuk: Resultaten en Evaluatie
%%=============================================================================

\chapter{Resultaten en Evaluatie}%
\label{ch:resultaten-en-evaluatie}

Dit hoofdstuk presenteert en interpreteert de meetresultaten van de proof-of-concept uit
Hoofdstuk~\ref{ch:proof-of-concept}. Voor elk van de drie afstandscategorieën uit
Hoofdstuk~\ref{ch:requirementsanalyse} (0--15~km, 15--30~km en 30+~km) werden honderd
onafhankelijke trials uitgevoerd met \href{code/scripts/Run-DistanceTrials.ps1}{\texttt{Run-DistanceTrials.ps1}},
waarbij elke trial een eigen, willekeurige weers- en DFS-realisatie kreeg
(Sectie~\ref{sec:sim-parameters}). Dat levert in totaal 300 onafhankelijke metingen op,
waarmee de architectuur niet enkel op een enkel gunstig moment, maar over een
representatieve steekproef van mogelijke kanaalomstandigheden geëvalueerd kan worden.

\section{Methodologie van de metingen}%
\label{sec:res-methodologie}

Elke trial verliep volgens hetzelfde protocol: eerst kreeg de link-emulator via
\texttt{poc-run-trial.sh} een nieuwe, willekeurige combinatie van vertraging,
doorvoersnelheidsbeperking en Gilbert-Elliot-burstverlies opgelegd, gecentreerd rond de
basiswaarden van de gekozen afstandscategorie (Sectie~\ref{subsec:willekeurige-trial}).
Onafhankelijk daarvan kreeg elke trial een kans van 8\% op een bijkomende
DFS-onderbreking, ongeacht de afstandscategorie, zoals vereist. Vervolgens werd
gedurende een meetvenster van 30~seconden het effectieve packet loss tussen client en
server gemeten met \texttt{ping} (30 pakketten, 1 per seconde) -- een venster dat
ruimschoots volstaat om een eventuele DFS-onderbreking (uiterlijk na 18~seconden
volledig afgerond) mee op te vangen.

Deze batchmethodologie is specifiek gericht op REQ-03 (packet loss-tolerantie) en FR-06
(hervatten na onderbreking). De doorvoersnelheid (REQ-01) en de latency/jitter (REQ-02)
werden, zoals beschreven in Sectie~\ref{sec:sim-uitvoering}, afzonderlijk gevalideerd via
\texttt{Run-Tests.ps1} en vallen buiten de omvang van de 300 batchtrials in dit
hoofdstuk.

Om een zuiver antwoord te geven op de vraag of de architectuur onder \emph{normale}
kanaalomstandigheden aan REQ-03 voldoet, worden de resultaten per afstandscategorie
telkens in drie groepen opgesplitst: alle trials samen (het ruwe gemiddelde), enkel de
trials \emph{zonder} DFS-gebeurtenis (relevant voor REQ-03), en enkel de trials
\emph{met} een DFS-gebeurtenis (relevant voor FR-06). Deze opsplitsing werd berekend met
\href{code/scripts/Analyze-Trials.ps1}{\texttt{Analyze-Trials.ps1}} op basis van de
ruwe, per trial weggeschreven CSV-bestanden.

\section{Resultaten per afstandscategorie}%
\label{sec:res-resultaten}

Tabel~\ref{tab:resultaten-per-afstand} vat de statistische verdeling van het gemeten
packet loss samen voor elk van de drie afstandscategorieën, telkens opgesplitst volgens
Sectie~\ref{sec:res-methodologie}. Figuur~\ref{fig:packetloss_per_afstand} toont
dezelfde gegevens grafisch.

\begin{table}[h]
  \centering
  \small
  \begin{tabular}{p{2.0cm}p{2.6cm}p{1.3cm}p{1.9cm}p{1.9cm}p{1.7cm}p{1.9cm}}
    \toprule
    \textbf{Afstand} & \textbf{Groep} & \textbf{$n$} & \textbf{Gem.} & \textbf{Std.dev.} & \textbf{Min/Max} & \textbf{$>$REQ-03} \\
    \midrule
    0--15~km & Alle trials        & 100 & 2{,}43\%  & 8{,}31\%  & 0{,}00 / 36{,}67\% & 15/100 \\
    0--15~km & Zonder DFS         & 94  & 0{,}35\%  & 1{,}14\%  & 0{,}00 / 6{,}67\%  & 9/94   \\
    0--15~km & Met DFS            & 6   & 35{,}00\% & 1{,}67\%  & 33{,}33 / 36{,}67\% & 6/6   \\
    \midrule
    15--30~km & Alle trials        & 100 & 7{,}37\%  & 8{,}87\%  & 0{,}00 / 36{,}67\% & 78/100 \\
    15--30~km & Zonder DFS         & 93  & 5{,}27\%  & 4{,}65\%  & 0{,}00 / 20{,}00\% & 71/93  \\
    15--30~km & Met DFS            & 7   & 35{,}24\% & 1{,}65\%  & 33{,}33 / 36{,}67\% & 7/7    \\
    \midrule
    30+~km & Alle trials        & 100 & 21{,}47\% & 10{,}85\% & 0{,}00 / 53{,}33\% & 99/100 \\
    30+~km & Zonder DFS         & 96  & 20{,}38\% & 9{,}62\%  & 0{,}00 / 46{,}67\% & 95/96  \\
    30+~km & Met DFS            & 4   & 47{,}50\% & 3{,}63\%  & 43{,}33 / 53{,}33\% & 4/4    \\
    \bottomrule
  \end{tabular}
  \caption{Statistische verdeling van het gemeten packet loss per afstandscategorie
  (100 trials elk), opgesplitst naar de aan- of afwezigheid van een DFS-gebeurtenis
  tijdens de trial. "$>$REQ-03" geeft het aantal trials weer dat de norm van maximaal
  2\% packet loss overschrijdt.}
  \label{tab:resultaten-per-afstand}
\end{table}

\begin{figure}[h]
  \centering
  \includegraphics[width=\textwidth]{packetloss_per_afstand.png}
  \caption{Gemiddeld packet loss per afstandscategorie, opgesplitst naar trials zonder
  (groen, REQ-03) en met (rood, FR-06) DFS-gebeurtenis. De stippellijn markeert de
  REQ-03-norm van 2\%. Foutbalken tonen de standaardafwijking.}
  \label{fig:packetloss_per_afstand}
\end{figure}

\section{Interpretatie}%
\label{sec:res-interpretatie}

De drie afstandscategorieën vertonen een duidelijk, oplopend patroon dat niet toevallig
is, maar rechtstreeks voortvloeit uit de progressief zwaardere basisparameters die in
Sectie~\ref{sec:sim-parameters} per categorie werden vastgelegd. Toch verschilt de
\emph{oorzaak} van eventuele REQ-03-overschrijding sterk per categorie, wat enkel
zichtbaar wordt dankzij de DFS-opsplitsing.

\paragraph{0--15~km: REQ-03 wordt gehaald onder normale werking.} Zonder DFS-gebeurtenis
bedraagt het gemiddelde packet loss slechts 0{,}35\%, ruim onder de norm van 2\%. Van de
94 trials zonder DFS overschrijden er slechts 9 de norm, telkens met een beperkte
overschrijding (maximaal 6{,}67\%). De overige 6 trials met een REQ-03-overschrijding
gaan uitsluitend terug op een DFS-gebeurtenis: bij die trials ligt het verlies steevast
rond 33--37\%, een waarde die rechtstreeks te verklaren is uit een volledige
onderbreking van ongeveer 10~seconden binnen een meetvenster van 30~seconden
($10/30 \approx 33\%$). Voor deze afstandscategorie is REQ-03 dus een correcte
weerspiegeling van een goed presterend kanaal: de architectuur voldoet aan de eis onder
normale omstandigheden, en de enige uitzonderingen zijn een rechtstreeks, verwacht
gevolg van een verplichte regelgevende gebeurtenis (Sectie~\ref{subsec:dfs}) -- precies
het scenario waarvoor FR-06 bedoeld is, niet een tekortkoming van de link zelf.

\paragraph{15--30~km: REQ-03 wordt structureel overschreden, ook zonder DFS.} Hier
bedraagt het gemiddelde packet loss zonder DFS al 5{,}27\%, meer dan het dubbele van de
norm, en overschrijden 71 van de 93 niet-DFS-trials (76{,}3\%) de norm. In tegenstelling
tot de 0--15~km-categorie is de REQ-03-overschrijding hier dus \emph{niet} voornamelijk
een DFS-effect, maar een gevolg van de courante, alledaagse kanaalkwaliteit op deze
afstand. De DFS-trials (gemiddeld 35{,}24\%) volgen wel nog steeds hetzelfde patroon als
bij 0--15~km, wat aantoont dat het herstelgedrag na een onderbreking (FR-06) op deze
afstand identiek robuust blijft; het is de basiskwaliteit van de verbinding zelf die
tekortschiet.

\paragraph{30+~km: de architectuur voldoet niet aan REQ-03.} Zonder DFS bedraagt het
gemiddelde packet loss 20{,}38\%, en overschrijden 95 van de 96 niet-DFS-trials de norm.
Op deze afstand is de REQ-03-overschrijding niet langer een grensgeval, maar
vrijwel universeel. Dit bevestigt kwantitatief de kwalitatieve inschatting uit
Sectie~\ref{sec:req-tussentijdse-conclusie}, waar 30+~km reeds werd omschreven als
liggend buiten het betrouwbare toepassingsgebied voor permanente bestandstoegang.

\paragraph{FR-06 (herstel na DFS-onderbreking) wordt op elke afstand consistent
gevalideerd.} Opvallend is dat het gemiddelde verlies bij DFS-trials weliswaar toeneemt
met de afstand (35{,}00\% → 35{,}24\% → 47{,}50\%), maar dat het patroon -- een pak van
het meetvenster volledig onbereikbaar, gevolgd door herstel binnen dezelfde trial --
op alle drie de afstanden aanwezig blijft. Dit toont aan dat de architectuur
(SFTP over TCP, in combinatie met de netem-gebaseerde DFS-simulatie) consequent in staat
is om van een verplichte kanaalonderbreking te herstellen, onafhankelijk van de
onderliggende afstandscategorie.

\section{Toetsing aan de requirements}%
\label{sec:res-toetsing}

Tabel~\ref{tab:res-toetsing} vat de toetsing van de statistische batchresultaten aan de
relevante requirements uit Hoofdstuk~\ref{ch:requirementsanalyse} samen.

\begin{table}[h]
  \centering
  \small
  \begin{tabular}{p{1.6cm}p{9.0cm}p{2.4cm}}
    \toprule
    \textbf{Vereiste} & \textbf{Toetsing} & \textbf{Resultaat} \\
    \midrule
    REQ-03 (0--15~km) & Packet loss $\le 2\%$ onder normale werking & \textbf{Voldaan} (0{,}35\%) \\
    REQ-03 (15--30~km) & Packet loss $\le 2\%$ onder normale werking & \textbf{Niet voldaan} (5{,}27\%) \\
    REQ-03 (30+~km) & Packet loss $\le 2\%$ onder normale werking & \textbf{Niet voldaan} (20{,}38\%) \\
    FR-06 (alle afstanden) & Herstel van SFTP-overdracht na een DFS-onderbreking & \textbf{Voldaan} op alle drie de afstanden \\
    \bottomrule
  \end{tabular}
  \caption{Toetsing van de statistische batchresultaten aan REQ-03 en FR-06.}
  \label{tab:res-toetsing}
\end{table}

\section{Besluit van dit hoofdstuk}%
\label{sec:res-besluit}

De 300 onafhankelijke trials, verdeeld over drie afstandscategorieën en met telkens een
eigen willekeurige weers- en DFS-realisatie, leveren een statistisch onderbouwd en
consistent beeld op. Binnen het optimale bereik (0--15~km) presteert de architectuur
ruim binnen de gestelde norm voor packet loss, met REQ-03-overschrijdingen die vrijwel
uitsluitend het rechtstreekse en verwachte gevolg zijn van een verplichte
DFS-onderbreking. Vanaf de horizonlimiet (15--30~km) schiet de courante kanaalkwaliteit
zelf reeds structureel tekort, en voorbij het als onbetrouwbaar bestempelde bereik
(30+~km) is die tekortkoming vrijwel universeel. Onafhankelijk van de afstand toont de
architectuur zich wel steeds in staat om een SFTP-overdracht na een DFS-onderbreking te
hervatten (FR-06), wat de robuustheid van de gekozen combinatie van SFTP en de
TCP-transportlaag bevestigt.

Deze resultaten geven, in combinatie met de eerder afzonderlijk gevalideerde
doorvoersnelheid en latency (REQ-01, REQ-02, Sectie~\ref{sec:sim-uitvoering}), een
genuanceerd antwoord op de centrale onderzoeksvraag: de voorgestelde architectuur is
geschikt voor toegang tot een private fileserver binnen het optimale bereik, maar niet
zonder meer voor de verder gelegen afstandscategorieën. De implicaties hiervan voor de
onderzoeksvraag en de aanbevelingen voor vervolgonderzoek worden verder uitgewerkt in
Hoofdstuk~\ref{ch:conclusie}.